% TOP_PROBLEM: {"id": 146, "title": "Irreducibility of Random {0,1} Polynomials", "category_id": 1, "category": {"id": 1, "name": "number_theory", "display_name": "Number Theory", "description": "Properties of integers, prime numbers, Diophantine equations.", "slug": "number-theory", "order_index": 1, "created_at": "2026-07-31T15:26:25.670Z"}, "status": "open"}
% FIRST_POSTED: null
% ENTRY_KIND: statement_only
% Imported from: batch03.tex; UnsolvedMath ID 146
\documentclass[11pt]{article}
\usepackage[margin=25mm]{geometry}
\usepackage{fontspec}
\setmainfont{Latin Modern Roman}
\usepackage{amsmath,amssymb,mathtools}
\usepackage{unicode-math}
\setmathfont{Latin Modern Math}
\usepackage{xurl,hyperref}
\hypersetup{colorlinks=true,urlcolor=blue}
\setcounter{secnumdepth}{0}
\setlength{\parindent}{0pt}
\setlength{\parskip}{5pt}
\setlength{\emergencystretch}{3em}
\newcommand{\sref}[2]{\hyperlink{source:#1:#2}{[S#2]}}
\newcommand{\cataloguescope}[1]{\paragraph{Recorded target.}#1}
\begin{document}
% UnsolvedMath ID 146; source catalogue code GREEN-058
\setcounter{section}{145}
\section{Irreducibility of Random \{0,1\} Polynomials}\label{146}
{\small\sffamily UnsolvedMath ID 146 · Number Theory}
\subsection{Definitions and mathematical statement}

\paragraph{Shared notation.}$\mathbb N=\{1,2,\ldots\}$, and $\mathbb Z,\mathbb Q,\mathbb R,\mathbb C$ denote the integers, rationals, reals and complex numbers. Any other conventions are specified locally.


For each integer $n\geq2$, let $\xi_1,\ldots,\xi_{n-1}$ be independent random variables with
$\Pr(\xi_j=0)=\Pr(\xi_j=1)=1/2$, and set
\[
 F_n(T)=T^n+\sum_{j=1}^{n-1}\xi_jT^j+1\in\mathbb Z[T].
\]
This is the uniform distribution on degree-$n$ polynomials with coefficients in $\{0,1\}$ and nonzero constant coefficient; the leading and constant coefficients are necessarily $1$. Irreducible means having no factorization into two polynomials of positive degree in $\mathbb Q[T]$.
\paragraph{Statement recorded in the supplied source.}
Does
\[
 \lim_{n\to\infty}\Pr\bigl(F_n\text{ is irreducible over }\mathbb Q\bigr)=1?
\]
Here ``almost surely'' means probability tending to one as the degree tends to infinity, rather than an assertion about every term of one infinite random sequence. \sref{146}{2}
\subsection{Short English statement}
Is a uniformly random binary-coefficient polynomial with constant coefficient one irreducible with probability tending to one as its degree grows?
\subsection{Sources}
\begin{description}
\item[\hypertarget{source:146:1}{S1}] ULAM.AI and the UnsolvedMath Contributors, \emph{UnsolvedMath: A Curated Collection of Open Mathematics Problems}, record 146 (GREEN-058). CC BY 4.0. \url{https://huggingface.co/datasets/ulamai/UnsolvedMath}.
\item[\hypertarget{source:146:2}{S2}] Lior Bary-Soroker, Dimitris Koukoulopoulos and Gady Kozma, \emph{Irreducibility of random polynomials: general measures}, arXiv:2007.14567v3 (2023), §1, discussion of binary coefficients and Theorem 1 with coefficient set $\{0,1\}$. \url{https://arxiv.org/abs/2007.14567}.
\item[\hypertarget{source:146:3}{S3}] Ben Green, \emph{100 Open Problems}, maintained author list, Problem 93 and its comments, accessed 7 October 2026. \url{https://people.maths.ox.ac.uk/greenbj/papers/open-problems.pdf}.
\end{description}
\label{end:146}
\end{document}
